% ------------------------------------------------------------
% English working translation layer.
% Source file: papers_60_90/66_La_serie_exceptionnelle_de_groupes_de_Lie_II/66_La_serie_exceptionnelle_de_groupes_de_Lie_II.tex
% Modern title: The exceptional series of Lie groups II
% ------------------------------------------------------------

\documentclass[11pt,a4paper]{article}
\usepackage[utf8]{inputenc}
\usepackage[T1]{fontenc}
\usepackage[english]{babel}
\usepackage{amsmath,amssymb,amsthm}
\usepackage{amsfonts}
\usepackage{mathtools}
\usepackage{geometry}
\usepackage{hyperref}
\usepackage{enumitem}
\usepackage{booktabs}
\usepackage{array}
\usepackage{mathrsfs}
\usepackage{calrsfs}
\usepackage{xcolor}
\usepackage{microtype}

\geometry{a4paper, top=2.5cm, bottom=2.5cm, left=2.5cm, right=2.5cm}

\theoremstyle{plain}
\newtheorem{theorem}{Th\'eor\`eme}
\newtheorem{lemma}[theorem]{Lemma}
\newtheorem{proposition}[theorem]{Proposition}
\newtheorem{corollary}[theorem]{Corollary}
\newtheorem{conjecture}[theorem]{Conjecture}

\theoremstyle{definition}
\newtheorem{definition}[theorem]{D\'efinition}
\newtheorem{example}[theorem]{Example}
\newtheorem{remark}[theorem]{Remark}

\DeclareMathOperator{\Tr}{Tr}
\DeclareMathOperator{\ad}{ad}
\DeclareMathOperator{\Hom}{Hom}
\DeclareMathOperator{\rank}{rang}
\newcommand{\rang}{\rank}

\begin{document}
\begin{center}\small\emph{Literal English translation working draft. Formulae, numbering, and TeX structure are preserved; this is not yet a modern recast.}\end{center}
\vspace{0.5em}


\noindent\textbf{Algebras of Lie/Lie Algebras}

\begin{center}
\Large\textbf{The exceptional series of Lie groups II}
\end{center}

\begin{center}
\textbf{Pierre DELIGNE and Ronald of MAN}
\end{center}

\noindent P.~D.: Institute for Advanced Study, School of Mathematics, Princeton, NJ 08540, USA;\\
E-mail: Deligne@math.ias.edu

\medskip
\noindent R.~of M.: Fac.\ Wisk.\ and Inf., Technical University Eindhoven,\\
PO Box 513, 5600 MB Eindhoven, The Netherlands.

\medskip
\noindent\textbf{Resume.} Numerologie of the groups exceptionnels and of leurs representations.

\medskip
\noindent\textbf{Abstract.} Numerology of exceptional Lie groups and their representations.

\bigskip

As in [1] and [2], we considerons in this note the groups $G$ of the serie exceptional
\[
\text{e},\ A_1,\ A_2,\ G_2,\ D_4,\ F_4,\ E_6,\ E_7,\ E_8.
\]
Between the group trivial and $A_1$, we intercalons en outre the super-group $\mathrm{SOSp}(1,2)$. More that the group, this which we interesse is each fois the category of ses representations. For $\mathrm{SOSp}(1,2)$, one considers the category of the super-representations such that the parite let given by $-1 \in \mathrm{Sp}(2) \subset \mathrm{SOSp}(1,2)$. The composante neutre $G^0$ of each group is the group adjoint, and the category of representations considered is each fois $\otimes$-generated by the representation adjointe.

The nombre of Coxeter dual $h^\vee$ of $G^0$ is, selon the type: $6/5$, $3/2$, $2$, $3$, $4$, $6$, $9$, $12$, $18$, $30$. For the group trivial, $h^\vee = 6/5$ by fiat. We considerons $G$ as dependant of the parametre $h^\vee$. Parfois, it is more commode of utiliser $\mu := 6/h^\vee$, that we preferons maintenant a $\lambda$ utilise in [1] and [2]: $\mu = -\lambda$. Values of $\mu$: $5, 4, 3, 2, 3/2, 1, 2/3, 1/2, 1/3, 1/5$.

For $G^0$ of rank $r$, considerons the weight dominants $\theta_1,\dots,\theta_r$ suivants. For $E_8$, we we arretons a $\theta_7$. Notation: $(a_1,\dots,a_r)$ for $\sum a_i \theta_i$, the weight fondamentaux $\omega_i$ being pris in the ordre of the tables of Bourbaki. For $G^0$ of rank $\geq 1$ (i.e.\ $G \not= \text{e}$), $\theta_1$ is the more grande racine, corresponding a the representation adjointe. This defines the $\theta_i$ for $G$ of rank $\leq 1$, i.e.\ $G^0 = \text{e}$, $\mathrm{SOSp}(1,2)$ or $A_1$. Ensuite, one prend:
\begin{itemize}[leftmargin=2em]
\item[] $A_2$: $(1,1)$, $(0,3)$;
\item[] $G_2$: $(0,1)$, $(3,0)$;
\item[] $D_4$: $(0,1,0,0)$, $(1,0,1,1)$, $(2,0,0,0)$, $(2,0,2,0)$;
\item[] $F_4$: $(1,0,0,0)$, $(0,1,0,0)$, $(0,0,0,2)$, $(0,0,2,0)$;
\item[] $E_6$: $(0,1,0,0,0,0)$, $(0,0,0,1,0,0)$, $(1,0,0,0,0,1)$, $(0,0,1,0,1,0)$, $(0,0,0,0,1,1)$, $(1,0,0,0,2,0)$;
\item[] $E_7$: the weight fondamentaux $\omega_1$, $\omega_3$, $\omega_6$, $\omega_4$ and $(0,1,0,0,0,0,1)$, $(0,1,0,0,1,0,0)$, $(0,0,0,0,1,0,1)$;
\item[] $E_8$: the weight fondamentaux $\omega_8$, $\omega_7$, $\omega_1$, $\omega_6$, $\omega_2$, $\omega_5$, $\omega_3$.
\end{itemize}
For $A_2$, $D_4$ and $E_6$, it serait indifferent of remplacer these $\theta_i$ by leurs images by a same automorphism of the diagram of Dynkin.

For $G$ connected, we shall denote $\mathcal{R}(p_1,\dots,p_r)$ the representation irreducible of $G$ of weight dominant $\sum p_i \theta_i$. For $G$ disconnexe, this weight must be decore, as in [1], for define $\mathcal{R}(p_1,\dots,p_r)$. The decoration requise is the suivante. For $A_1$: the signe $(-1)^{p_1}$ si $p_2 = 0$, $0$ otherwise. For $D_4$: $+$, $0+$, or $0-$. For $E_8$: $+$ or $0-$.

Let $V_i$ the representation $\mathcal{R}(p_1,\dots,p_r)$ for $p_i = \delta_i^j$. The representation $V_1$, defined for $G$ of rank $\geq 1$, is the representation adjointe. The representations $V_i$ ($1 \leq i \leq 7$) are the representations $\mathcal{J}$, $\mathcal{X}_2$, $\mathcal{Y}_2^*$, $\mathcal{X}_3$, $\mathcal{O}^*$, $\mathcal{X}_4$, $\mathcal{F}^*$ of [1]. The representation $\mathcal{R}(p_1,\dots,p_r)$ is caracterisee by the properties of admettre the weight dominant $\sum p_i \theta_i$ and of be a constituant of the tensor product of the $V_i^{\otimes p_i}$.

For $s \leq r$, one ecrira again $\mathcal{R}(p_1,\dots,p_s)$ for $\mathcal{R}(p_1,\dots,p_s,0,\dots,0)$.

Notre but is of decrire, a $p_1,\dots,p_7$ fixes, of the similitudes between the representations $\mathcal{R}(p_1,\dots,p_s)$ of the groups of the serie exceptional of rank $\geq s$. Analogue in the serie $A$: a sequence decroissante $\ell_1 \geq \ell_2 \geq \cdots \geq \ell_k \geq 0$ of integers being datum, prendre the representation corresponding of $\mathrm{GL}(n)$, for each $n \geq k$.

\bigskip

Let $G$ a group simple deploye epingle, $\mathfrak{g}$ son algebra of Lie and $\mathfrak{h}$ the algebra of Lie of the tore deploye epinglant $T$. One identifie the group of the weight $X := \Hom(T,\mathbb{G}_\mathrm{m})$ a under-group of the dual $\mathfrak{h}^\vee$ of $\mathfrak{h}$. The form bilineaire canonical $(\, , \,)$ is the form bilineaire on $\mathfrak{h}^\vee$ inverse of the form of Killing $\Tr(\ad x \cdot \ad y)$ on $\mathfrak{h}$. For $G$ a super-group classique, by example $\mathrm{SOSp}(1,2)$: utiliser the super-trace. Si $\alpha_0$ is the more grande racine, a of the definitions of the nombre of Coxeter dual $h^\vee$ is that $(\alpha_0,\alpha_0) = 1/h^\vee$.

\`A each weight dominant $\lambda$, we shall attacher a systeme finite of nombres rationnels, chacun pris with a multiplicite $m(q) \in \mathbb{Z}$. En of other termes: a mesure $L_G(\lambda)$ on $\mathbb{Q}$, combinaison $\mathbb{Z}$-lineaire $\sum m(q)\delta[q]$ of mesures of Dirac. Let $[R^+]$ the mesure on $X$ sum of the mesures of Dirac $\delta[\alpha]$ for $\alpha$ a racine positive. Let $M(\lambda)$ the mesure on $\mathbb{Q}$ image of $[R^+]$ by $(2\lambda, \cdot): X \to \mathbb{Q}$. The ``12'' a for raison of be the assertion of integralite in the theorem~1. Si all the racines have the same longueur, or si $\alpha$ is a racine longue, $(\lambda,\alpha) = \lambda(H_\alpha)/2h^\vee$. Si $G$ admet of the racines of deux longueurs differentes and that the racine $\alpha$ is courte, $(\lambda,\alpha) = \bigl((\alpha,\alpha)/(\alpha_0,\alpha_0)\bigr) \cdot \lambda(H_\alpha)/2h^\vee$. The factor $(\alpha,\alpha)/(\alpha_0,\alpha_0)$ vaut $1/4$, $1/3$, $1/2$ for $\mathrm{SOSp}(1,2)$, $G_2$, $F_4$. Let $\rho$ the demi-sum $\frac{1}{2}\int \alpha \cdot [R^+]$ of the racines positives. For a super-group classique, by example $\mathrm{SOSp}(1,2)$, there is lieu in these definitions of remplacer $[R^+]$ by the sum on the racines positives $\sum m(\alpha)\delta[\alpha]$, with $m(\alpha) = 1$ (resp.\ $m(\alpha) = -1$) si the under-space radiciel corresponding is pair (resp.\ impair). We definissons
\begin{equation}\label{eq:1}
L_G(\lambda) := M(\lambda+\rho) - M(\rho).
\end{equation}

Si that ne cree not of ambiguite, one ecrira simplement $L(\lambda)$ for $L_G(\lambda)$. Si $L(\lambda) = \sum m(q)\delta[q]$, the formule of dimension of Weyl dit that the representation $V(\lambda)$ of weight dominant $\lambda$ is of dimension
\begin{equation}\label{eq:2}
\dim V(\lambda) = \prod_q m(q)^{-1}.
\end{equation}

Let $F_{2N}^{\mathfrak{g}}$ the function $\Tr((\ad x)^{2N})$ on $\mathfrak{g}$. Restreignons-the a $\mathfrak{h}$, and identifions $\mathfrak{h}$ a $\mathfrak{h}^\vee$ by the form bilineaire canonical. One obtains a function $F_{2N}^{\mathfrak{g}}$ on $\mathfrak{h}^\vee$. By construction,
\begin{equation}\label{eq:3}
F_{2N}^{\mathfrak{g}}(\lambda + \rho) - F_{2N}^{\mathfrak{g}}(\rho) = 2 \cdot 12^{-N} \sum m(q) q^{2N} L_G(\lambda).
\end{equation}
The factor $2$ is of the a this that the definition of $L_G(\lambda)$ part of $R^+$ and not of the set $R = R^+ \cup (-R^+)$ of all the racines. Appliquant the theorem 2, p.~267 of Duflo [3] to the weight $\lambda$ and $0$, it is possible of expliciter a element of the centre of the algebra enveloppante agissant on $V(\lambda)$ by $F_{2N}^{\mathfrak{g}}(\lambda+\rho) - F_{2N}^{\mathfrak{g}}(\rho)$. For $N = 1$, c'is the Casimir.

\smallskip

For $p_1,\dots,p_7 \in \mathbb{N}$, we construirons a combinaison lineaire formelle $L(p_1,\dots,p_7)$ of elements of the space $\mathcal{L}$ of the forms lineaires $x \mapsto Ax + B$ a coefficients integers. En of other termes: a combinaison $\mathbb{Z}$-lineaire $\sum m(A,B)\delta[Ax + B]$ of mesures of Dirac on $\mathbb{Q}$. The vertu of $L(p_1,\dots,p_7)$ is of interpoler the $L(\lambda)$, $\lambda$ weight dominant of $\mathcal{R}(p_1,\dots,p_7)$:

\begin{theorem}\label{thm:1}
For $p_1,\dots,p_7 \geq 0$ and for each group $G$ of the serie exceptional, si $p_i = 0$ for $i > r = \inf(\rang(G),7)$, and that $\lambda$ is a weight dominant of $\mathcal{R}(p_1,\dots,p_r)$, the mesure $L_G(\lambda)$ is the image of $L(p_1,\dots,p_7)$ by the map of $\mathcal{L}$ in $\mathbb{Q}$ << evaluation en $\mu$ >>: $x \mapsto x(\mu)$.
\end{theorem}

Si $L(p_1,\dots,p_7) = \sum m(A,B)\delta[Ax + B]$, $L(\lambda)$ is therefore the sum of the $m(A,B)\delta[A\mu + B]$ and (3) se recrit
\begin{equation}\label{eq:4}
F_{2N}^{\mathfrak{g}}(\lambda + \rho) - F_{2N}^{\mathfrak{g}}(\rho) = 2 \cdot 12^{-N} \sum m(A,B) (A\mu + B)^{2N}.
\end{equation}

Let $(\, , \,)$ the multiple of the form bilineaire canonical such that $(\alpha_0,\alpha_0) = 2$. En termes of $(\, , \,)$, (1), the theorem 1 and (4) se reformulent as suit. The difference between the images of the mesure $[R^+]$ by $(\lambda + \rho, \cdot)$ and by $(\rho, \cdot)$ is:
\begin{equation}\label{eq:5}
\langle \lambda + \rho, [R^+] \rangle - \langle \rho, [R^+] \rangle = \sum m(A,B)\delta[A + Bh^\vee/6],
\end{equation}
and
\begin{equation}\label{eq:6}
12^{-N} \bigl(F_{2N}^{\mathfrak{g}}(\lambda + \rho) - F_{2N}^{\mathfrak{g}}(\rho)\bigr) = \sum m(A,B)(A + Bh^\vee/6)^{2N}.
\end{equation}

of after the theorem 1, a $p_1,\dots,p_7$ fixes, the function $F_{2N}^{\mathfrak{g}}(\lambda + \rho) - F_{2N}^{\mathfrak{g}}(\rho)$ of $G$ is of degree $\leq 2N$ en the parametre $\mu = 6/h^\vee$ attached to $G$. For $N = 1$, i.e.\ for the scalaire by lequel Casimir agit on $\mathcal{R}(p_1,\dots,p_7)$, c'is same a function lineaire:

\begin{theorem}\label{thm:2}
With the notations precedentes, quels that let $p_1,\dots,p_7$, the sum $\sum m(A,B)(Ax + B)^2$ is lineaire en $x$, i.e.
\begin{equation}\label{eq:7}
\sum m(A,B) A^2 = 0.
\end{equation}
\end{theorem}

For $r \leq 7$, we ecrirons $L(p_1,\dots,p_r)$ for $L(p_1,\dots,p_r,0,\dots,0)$ ($7-r$ zeros finaux). Fixons $r \leq 7$ and $G$ in the serie exceptional of rank $\geq r$. A $G$ correspond $\mu \in \mathbb{Q}$.

\begin{theorem}\label{thm:3}
Si $p_1,\dots,p_r \geq 0$, the form lineaire zero n'is not in the support of $L(p_1,\dots,p_r)$, and the masse totale of the forms lineaires $Ax + B$ such that $A\mu + B = 0$ is zero. The scalaire
\[
\sum_{A\mu + B = 0} (A\mu + B) \cdot m(A,B)
\]
is a integer. C'is the longueur of the restriction a $G^0$ of $\mathcal{R}(p_1,\dots,p_r)$. of after (2), $\dim \mathcal{R}(p_1,\dots,p_r)$ is therefore the value en $\mu$ of the function rational $\prod (Ax + B)^{-m(A,B)}$.

More precisely, si $r \leq \rang(G)$ and that $A\mu + B = 0$, the multiplicite $m(A,B)$ of $Ax + B$ in $L(p_1,\dots,p_r)$ is zero sauf in the cas suivants. Si $G = A_2$ and that $p_2 > 0$: $1$ for $-2x + 4$ and $-1$ for $-2x + 2$; si $G = D_4, E_8$ and that $p_3 > 0$, $p_4 = 0$ or that $p_3 = 0$, $p_4 > 0$: $1$ for $-3x + 3$ and $-1$ for $-x + 1$; si $G = F_4, E_8$ and that $p_3, p_4 > 0$: $1$ for $-3x + 3$ and $-2x + 2$, $-2$ for $-x + 1$; $G = E_6$ and $p_5$ or $p_6 > 0$: $1$ for $-4x + 2$ and $-1$ for $-2x + 1$.
\end{theorem}

The $\mathcal{R}(p_1,\dots,p_7)$ for $p_1 + 2p_2 + 2p_3 + 3p_4 + 3p_5 + 3p_6 + 4p_7 \leq 4$ are the representations $1$, $\mathcal{J}$, $\mathcal{A}$, $\mathcal{C}$, $\mathcal{O}^*$, $\mathcal{D}$, $\mathcal{E}$, $\mathcal{F}$, $\mathcal{F}^*$, $\mathcal{G}$, $\mathcal{H}$, $\mathcal{I}$, $\mathcal{J}$, $\mathcal{X}_i$ ($i \leq 4$), $\mathcal{Y}_i$ ($i \leq 4$) and $\mathcal{Y}_i^*$ of [1], which se lisent on $E_7$ or $E_8$. The theorem 3 redonne ainsi a part of the formules of dimension of [1].

\smallskip

Let $\mathcal{V}(p_1,\dots,p_7)$ the convolution
\begin{equation}\label{eq:8}
\mathcal{V}(p_1,\dots,p_7) := L(p_1,\dots,p_7) * (\delta[x] - \delta[x-1]).
\end{equation}
Si $L(p_1,\dots,p_7) = \sum m(A,B)\delta[Ax + B]$ and that $\mathcal{V}(p_1,\dots,p_7) = \sum n(A,B)\delta[Ax + B]$, the functions a support finite $m(A,B)$ and $n(A,B)$ se determinent the a the other by
\begin{equation}\label{eq:9}
n(A,B) = m(A,B) - m(A+1,B),
\end{equation}
\begin{equation}\label{eq:10}
m(A,B) = \sum_{A' \geq A} n(A',B).
\end{equation}
In (10), the sum is on $A'$; the factor $(A' \geq A)$ vaut $1$ si $A' \geq A$ and $0$ otherwise.

We prendrons (10) for definition of $L(p_1,\dots,p_7)$, $\mathcal{V}(p_1,\dots,p_7)$ being defined by specialisation a partir of a combinaison $\mathbb{Z}$-lineaire $\mathbf{V}$ of mesures of Dirac on the space of the maps of $\mathbb{Z}^7$ in $\mathbb{Q}$ of the form
\begin{equation}\label{eq:11}
(B; n; a_1,\dots,a_7) \mapsto (2\textstyle\sum a_i p_i + n)x + B.
\end{equation}
The mesure $\mathbf{V}$ admet a description simple en termes of the systeme of racines of type $E_8$.

Ecrivons $(B; n; a_1,\dots,a_7)$ for the function (11), $\delta[B; n; a_1,\dots,a_7]$ for the mesure of Dirac localisee en this function, and $\{\varepsilon_1,\dots,\varepsilon_8\}$ for the racine $2\delta - \alpha_3$ of $E_8$. With these notations, the definition of $\mathbf{V}$ is the suivante.

\smallskip

(A) Each racine positive $\{c, e, g, h, f, d, b, a\}$ fixe by $s_4$ contribue a $\mathbf{V}$
\[
\delta[B+h+n-1; a,b,c,d,e,f,g] - \delta[B+h; n-1; a,b,c,d,e,f,g]
\]
avcc $5h + n = a + b + c + d + e + f + g + h$.

\smallskip

(B) Each paire of racines positives $\{c, e, g, h', f, d, b, a\}$, $\{c, e, g, h'', f, d, b, a\}$ permutees by $s_4$, with $h'' < h'$ (one has then $h'' = h' - 1$) contribue a $\mathbf{V}$
\[
\delta[B+h'+n'-1; a,b,c,d,e,f,g] - \delta[B+h''; n''-1; a,b,c,d,e,f,g]
\]
with $5h' + n' = a + b + c + d + e + f + g + h'$ and $5h'' + n'' = a + b + c + d + e + f + g + h''$.

\smallskip

(C) $\mathbf{V}$ is the sum of these contributions and of the combinaison $\mathbb{Z}$-lineaire of $\delta[B; n; 0,\dots,0]$ such that for each $B$ and $n$, the masse totale of the $(B; n; a_i,\dots)$ let zero.

\begin{definition}
(i) $\mathcal{V}(p_1,\dots,p_7)$ is the image of $\mathbf{V}$ by the map of evaluation en $(p_1,\dots,p_7)$:
\[
(B; n; a_1,\dots,a_7) \mapsto (2\textstyle\sum a_i p_i + n)x + B \in \mathbb{Q}.
\]

(ii) It follows of (C) ci-dessus that for each $B_0$, the masse for $\mathcal{V}(p_1,\dots,p_7)$ of the $Ax + B$ with $B = B_0$ is zero. One defines $L(p_1,\dots,p_7)$ by (8), i.e.\ by (10).
\end{definition}

It is immediat on the definition of $\mathbf{V}$ that

\begin{proposition}\label{prop:1}
\textup{(i)} On the support of $\mathbf{V}$, $B \in \mathbb{Z}$ and $n \geq 0$.

\textup{(ii)} The projections of $\mathbf{V}$ by
\begin{itemize}
\item[(a)] $(B; n; a_1,\dots,a_7) \mapsto (B; n) \in \mathbb{Z}^2$,
\item[(b)] $(B; n; a_1,\dots,a_7) \mapsto (a_1,\dots,a_7) \in \mathbb{Z}^7$,
\end{itemize}
are nulles.
\end{proposition}

The nullite (ii a) equivaut a $\mathcal{V}(0,\dots,0) = 0$, i.e.\ a $L(0,\dots,0) = 0$. This is compatible to the fait that $\mathcal{R}(0,\dots,0)$ is, for each $G$, the representation trivial, of weight dominant $0$, and that $L_G(0) = 0$.

Compte tenu of (1) and of this that the masse totale of $\mathbf{V}$ is zero, the nullite (ii b) equivaut a this that, for every $N$, as function of $p_1,\dots,p_7$,
\[
\sum (2\textstyle\sum a_i p_i + B)^{2N} L(p_1,\dots,p_7)
\]
is of degree $\leq 2N$. This is compatible, via (4), to the fait that for every Lie group simple $G$, $F_{2N}^{\mathfrak{g}}(\lambda)$ is a function of degree $\leq 2N$ en the weight $\lambda$.

Posons $\Lambda = 2\sum a_i \theta_i + \rho$. The theorem 2 can se reformuler
\begin{equation}\label{eq:12}
\int_\Lambda (\Lambda + x - 1)(\Lambda + 1) \cdot \mathbf{V} = 0,
\end{equation}
identiquement en the indeterminees $p_i$.

\bigskip

\textbf{Some indications on the proof of the theorem 1.} --- Let $G$ a group of the serie exceptional and $r = \inf(\rang(G),7)$. Sauf for $E_8$ and the racine simple $\alpha_4$, aucune racine positive n'is orthogonale a all the $\theta_i$ ($1 \leq i \leq r$). One sets $[R^G] := [R^+]$ si $G \not= E_8$ and $[R^G] = [R^+] - \delta[\alpha_4]$ si $G = E_8$. Definissons the part mobile of $L_G(\sum p_i \theta_i)$ by
\begin{equation}\label{eq:13}
L_G^{\textup{mob}}(\textstyle\sum p_i \theta_i) = \langle M(\sum p_i \theta_i + \rho), [R^G] \rangle.
\end{equation}
The definition (1) se recrit $L_G(\sum p_i \theta_i) = L_G^{\textup{mob}}(\sum p_i \theta_i) - L_G^{\textup{mob}}(0)$.

The theorem 1 for $G$ n'utilise of $\mathbf{V}$ that son image $\mathbf{V}_G$ by
\begin{equation}\label{eq:14}
(B; n; a_1,\dots,a_7) \mapsto (2\textstyle\sum a_i p_i + n + B; a_1,\dots,a_r).
\end{equation}
Decomposons $\mathbf{V}_G$ en $\mathbf{V}_G^+ - \mathbf{V}_G^-$, the part mobile $\mathbf{V}_G^+$ (resp.\ fixe $\mathbf{V}_G^-$) ayant son support the or $(a_1,\dots,a_r) \not= 0$ (resp.\ $= 0$). The theorem 1 resulte of son analogue for the parties mobiles of $L_G(\sum p_i \theta_i)$ and $\mathbf{V}_G$: by the proposition 1 (ii a), the part fixe se recupere en faisant $x = 0$. For $(a_1,\dots,a_r) \not= 0$, let $R^{(a_1,\dots,a_r)}$ the set of the racines positives $\alpha$ such that $(\alpha, \theta_i) = a_i$ for $1 \leq i \leq r$. The form $(\, , \,)$ is as en (5)--(6). For $G \not= E_8$, each $R^{(a_1,\dots,a_r)}$ is reduit a element or is vide, and leur reunion is $R^+$. For $G = E_8$, each $R^{(a_1,\dots,a_r)}$ is a orbite of $\{s_3, s_4\}$ or is vide. Leur reunion is $R^+ - \{\alpha_4\}$. Let $\mathbf{V}_G^{(a_1,\dots,a_r)}$ the mesure on $\mathbb{Q}$ inverse image of $\mathbf{V}_G$ by $q \mapsto (B\cdot h^\vee/6 + n; a_1,\dots,a_r)$: the masse of $q$ is celle of the $(B; n; a_1,\dots,a_r,\dots)$ with $Bh^\vee/6 + n = q$. Let $[R^{(a_1,\dots,a_r)}]$ the sum of the $\delta[\alpha]$ for $\alpha \in R^{(a_1,\dots,a_r)}$. The theorem 1, for the parties mobiles of $L$ and $\mathbf{V}$, se decompose en the sum of the egalites
\begin{equation}\label{eq:15}
(p_i \cdot [R^{(a_1,\dots,a_r)}]) * (\delta[0] - \delta[-1]) = \mathbf{V}_G^{(a_1,\dots,a_r)}.
\end{equation}
Notre proof of (15) is by force brute. Ci-dessous, some properties of the systemes of racines en jeu which can aider a comprendre this which se passe.

Lorsque all the racines have the same longueur, $(\lambda, \alpha) = \lambda(H_\alpha)$ and $(\rho,\alpha)$ is the sum of the coefficients of $\alpha$, exprime as combinaison lineaire of racines simples. For $E_8$, $(\rho,\alpha) = 5$ and (15) resulte immediatement of the definition of $\mathbf{V}$.

Let $R$ the systeme of racines of type $A_1$, $A_2$, $D_4$, $E_6$ or $E_7$ and $r$ son rank. In the systeme of racines of type $E_8$, to the weight fondamentaux $\theta_1,\dots,\theta_r$ correspondent of the racines simples $\alpha_{S(1)},\dots,\alpha_{S(r)}$. Let $S_1$ the set complementaire of $(8-r)$ racines simples, and $W'$ the group of Weyl generated by the reflexions correspondantes. Let $\langle \theta_1,\dots,\theta_r \rangle$ the space vector generated by $\theta_1,\dots,\theta_r$. Other descriptions: the orthogonal of $\langle \theta_i \mid i \in S_1 \rangle$, the under-space fixe by $W'$. One can verify, cas by cas, that the racines of $E_8$ in $\langle \theta_1,\dots,\theta_r \rangle$ forment a systeme of racines of type $R$, and that si the one prend as racines positives of this systeme of racines celles which are positives for $E_8$, the weight fondamentaux of $E_8$ are the weight dominants $\theta_1,\dots,\theta_r$ of $R$.

For $R$ of the a of the types indiques, (15) resulte then of the assertions suivantes, that we have verifiees for $D_4$, $F_4$ and $E_7$:

(a) Let $T$ a orbite of $W'$ in the set of the racines of $E_8$. Si $T$ n'is not reduite a point fixe, and that for $t \in T$, $(t,\theta_i)$ (constant on $T$) is $> 0$ for a $\theta_i$ ($1 \leq i \leq r$), then the contribution of $T$ a $\mathbf{V}$ a image zero by (14).

(b) Notons $\rho_R$ and $h_R^\vee$ the sum of the weight fondamentaux, and the nombre of Coxeter dual, of $R$. Si $\alpha$ is in $\langle \theta_1,\dots,\theta_r \rangle$ is a racine positive, necessarily fixe by $s_4$, with the notations of the part (A) of the definition of $\mathbf{V}$, one has $(h^\vee/6)(\rho_R,\alpha) + n = (\rho,\alpha)$.

\smallskip

\textbf{Examples.} --- Verifions (a) for $E_7$. the orbite $T$ is ici a paire of racines as in the part (B) of the definition of $\mathbf{V}$. It s'agit of verify that $3h' + n' \geq 3h'' + n'' - 1$. Indeed, $h'' = h' - 1$ and $4h' + n' = 4h'' + n''$, i.e.\ $n'' = n' + 4$.

\smallskip

The theorem 1 given the dependance en $G$ of the $F_{2N}^{\mathfrak{g}}(\lambda + \rho) - F_{2N}^{\mathfrak{g}}(\rho)$, for certain $\lambda$.

\begin{proposition}\label{prop:2}
For $1 \leq N \leq 8$, $F_{2N}(\rho)$ is, as function of $G$, a function of $\mu = 6/h^\vee$ of the form $P_{2N}(\mu)/\mu(\mu+1)$, for a certain polynome $P_{2N}(X)$ of degree $2N$ verifiant $P_{2N}(X) = P_{2N}(-1-X)$.
\end{proposition}

It follows of Duflo [3] that $F_{2N}(\rho)$ is the value for $G$ (to the sens of [5]) of a combinaison lineaire of graphes trivalents a sommets orientes. For $N \leq 3$, of the arguments imites of [5] 6.8, 6.9 permettent of deduire the proposition of this presentation of $F_{2N}(\rho)$. To the-dela, we ne savons the verify that by force brute. Puisque we n'have that 10 groups $G$ a notre disposition, the proposition reste vraie, but is vide, for $N \geq 9$.

\bigskip

\textbf{Correction a [2].} --- The formule datum for $\dim(\mathcal{X}_4)$ is erronee. See [1] for the formule correcte.

\bigskip

\noindent\textbf{Note remise the 4 juin 1996, acceptee the 17 juin 1996.}

\bigskip

\begin{center}
\textbf{References bibliographiques}
\end{center}

\begin{enumerate}
\item[{[1]}] Cohen A.\ M.\ and of Man R., 1996. Computational evidence for Deligne's conjecture regarding exceptional Lie groups, \textit{C.\ R.\ Acad.\ Sci.\ Paris}, 322, serie I, p.~427--432.

\item[{[2]}] Deligne P., 1996. The exceptional series of Lie groups. \textit{C.\ R.\ Acad.\ Sci.\ Paris}, 322, serie I, p.~321--326.

\item[{[3]}] Duflo M., 1977. Operateurs differentiels bi-invariants on a Lie group. \textit{Ann.\ Sci.\ ENS}, 4e serie, 10, p.~265--283.

\item[{[4]}] van Leeuwen M.\ A.\ A., Cohen A.\ M.\ and Lisser B., 1992. \textit{LiE, a package for Lie group computations}, CAN, Amsterdam.

\item[{[5]}] Vogel P., 1995. Algebraic structures one modules of diagrams, preprint.
\end{enumerate}

\end{document}
